\section{Task Space, Kinematics, Trajectories, Precision, Continuity (Basic Robotics)}

\subsection{Task Coordinates}
Our task coordinates are the $x, y, z$ coordinates of the tip of the gripper when it is fully closed and the gripper's tilt angle,  $\theta_{tilt}$, relative to the table and the gripper's roll angle when facing the table $\theta_{roll}$. The tilt angle is measured relative to when the gripper is facing away from the table, which means that when the gripper is parallel with the table it is at 0 radians when facing along the table and $-\pi$ radians when facing towards it. The roll angle is relevant only when the gripper is facing down towards blocks on the table in preparation to pick or place them. The zero position of the gripper motor is when the line between its fingers is parallel with the x axis and the axis that the gripper rotates about being the negative world z axis in this position.

\subsection{Kinematics}
The tilt angle is simply computed in forward kinematics as the sum of the shoulder, elbow, and wrist tilt motors each multiplied by the sign that they move in (with the elbow being negative and the others being positive). Similarly, the roll angle is computed as the sum of the negative base angle and positive wrist roll motor angle. Therefore, the Jacobian for the last two task coordinates ($\theta_{tilt}$ and $\theta_{roll})$ are simply rows with ones or negative ones depending on direction on the joints that contribute to the angle. This results in the following matrix:

$$\begin{bmatrix}
0 & 1 & -1 & 1 & 0 \\
-1 & 0 & 0 & 0 & 1
\end{bmatrix}$$
\noindent We simply add these two rows to the bottom of the cartesian ($x,y,z$) Jacobian to compute those last two task entries.

\noindent The inverse kinematics is computed via the Newton-Raphson algorithm. For the Jacobian, we only use the first four task coordinates -- the cartesian coordinates and gripper tilt -- as the task vector and ignore the Jacobian's last row. This is because the gripper roll task is only useful when the gripper is pointing straight down towards the table and its degree of freedom corresponds directly to the gripper roll motor, so we can simply solve for it after the inverse kinematics as $q_{gripper\_roll} = \theta_{roll} + q_{base}$ based on the desired roll angle and the base motor's state outputted by the Newton-Raphson.

\subsection{Trajectories}
The arm's motion is governed by a trajectory queue---an ordered sequence of waypoints, each specifying a target joint configuration with velocity, a movement duration, and optional hold delays before and after the motion. The Manipulator node interpolates between consecutive waypoints using fifth-order polynomial splines, which guarantee continuous position and velocity with zero boundary accelerations for smooth motion. The Brain sends trajectory commands with a type field---replace, append-front, or append-back---that controls how new waypoints merge with any in-progress motion. This allows the system to chain complex multi-step maneuvers (e.g., grasp, lift, separate, grip-check, approach, place) as a continuous trajectory without stopping between phases.

\subsection{Regularization}
To ensure our robot was stable in the real-workspace, we imposed joint limits, used a damped pseudoinverse, and "restricted" to ensure elbow-up solutions. We limited the wrist-tilt motor to be $\in [-\frac{5\pi}{4}, \frac{pi}{4}]$ to ensure we don't impose a position on the motor such that neither the wrist-bracket collides ($-\frac{5\pi}{4}$) nor the camera mount collide with the upperarm link. We used a damped pseudo-inverse calculation to compute the inverse of the Jacobian via SVD, which ensures singular values below parameter gamma are regularised to avoid numerical instability near kinematic singularities. During early testing we noticed that sometimes the inverse kinematics solution would have an elbow-down configuration, which adds the risk of having the elbow motor colliding with the table. To account for this, we made our robot workspace slightly more conservative and made each joint guess to Newton-Raphson start from an elbow-up solution. From this point on, we didn't observe any elbow down ikin solutions.

\subsection{Gravity Compensation}

The gravity compensation model for the robot arm is formulated as a multi-body static-torque system. The model calculates the torques required at the shoulder, elbow, and wrist-tilt joints to counteract gravitational loads, assuming a vertical zero-configuration. The remaining joints (base, wrist-roll, and gripper) are assumed to be invariant to gravitational loading in the arm because they do not affect the extension of the arm and can not oppose the torque caused by gravity.

\noindent The model utilizes precomputed moment coefficients, $C_i$, for each joint. These coefficients aggregate the moment generated by a link's own mass ($m_i$) acting at its center of mass ($d_i$) and the moment generated by the total distal mass ($W_{distal}$) acting at the link's full length ($L_i$). This relationship is defined as:
\begin{equation}
    C_i = g (m_i d_i + W_{distal} L_i)
\end{equation}
where $g$ represents the acceleration due to gravity. An empirical tuning factor is included in the distal mass of the elbow to account for unmodeled hardware and mounting components, ensuring the system does not undershoot during compensation.

\noindent The required torques are calculated recursively from the end-effector back to the base. Let $\theta_1, \theta_{12},$ and $\theta_{123}$ represent the absolute angles of the shoulder, elbow, and wrist relative to the vertical axis. The gravitational torques ($\tau_g$) are determined by:
\begin{align}
    \tau_{g,w} &= C_3 \sin(\theta_{123}) \\
    \tau_{g,e} &= C_2 \sin(\theta_{12}) + \tau_{g,w} \\
    \tau_{g,s} &= C_1 \sin(\theta_1) + \tau_{g,e}
\end{align}

\noindent The final commanded motor torques are derived by negating the calculated gravitational torques and adjusting for the specific rotational direction ($d_i$) of each actuator:
\begin{equation}
    \tau_{motor,i} = -d_i \cdot \tau_{g,i}
\end{equation}

\subsection{Error Map}
We produced two calibration error maps to compensate for systematic forward-kinematics inaccuracies. The XY error map (Figure~\ref{fig:xy_error}) is a vector field showing the offset between commanded and actual tip positions across the workspace, measured during the XY calibration procedure described in Section~\ref{sec:xy_calib}. The Z error map (Figure~\ref{fig:z_error}) is a heatmap of measured table height versus radial distance from the base, captured during the Z calibration described in Section~\ref{sec:z_calib}. Both maps are stored as CSV files and loaded by the IK Solver at startup. During each IK solve, the solver interpolates from these maps to apply position-dependent corrections to the task coordinates before computing joint angles, compensating for systematic FK errors that vary with arm configuration.

\begin{figure}[H]
\centering
\begin{minipage}{.5\textwidth}
  \centering
  \includegraphics[width=1\linewidth]{final/imgs/xy_error_vector_field.png}
  \captionof{figure}{XY Error Map}
  \label{fig:xy_error}
\end{minipage}%
\begin{minipage}{.5\textwidth}
  \centering
  \includegraphics[width=.7\linewidth]{final/imgs/z_table_height.png}
  \captionof{figure}{Z Error Map}
  \label{fig:z_error}
\end{minipage}
\end{figure}

\subsection{Motion Detectors}
We implemented two simple motion detectors to determine contact collisions, gripper failures, and jerk movements. 

\noindent To detect collisions, we check for the difference between the commanded and actual joint positions, velocities, and effort and determine a collision occurred if any is above a threshold. We made the threshold be proportional to the speed of the robot which serves as a heuristic for incorporating inertia to our collision check. 

\noindent To determine a gripper failure, we check for the magnitude of the effort exerted by the gripper motor to achieve the commanded gripper position (fully closed) and signal a failure if the value is lower than a threshold. This is because if we are expecting a block to be gripped but it actually isn't then the gripper will be fully closed and the motor will not need to exert as much effort to try to fully get to the closed position.

\noindent As we continued developing our robot and the software for it, our required computations grew. We started to notice that as more compute was being spent, the higher the chance or the HEBI motors to lose connection and thus "drop" to a lower position. After connection would be gained, the motors would realize the commanded joint positions and move back quickly with a noticeable jerky movement. To determine jerk movements, we do the same logic as contact collision check but only on the difference between commanded and actual joint positions. Note that after we determine a jerk movement, we don't send the joint command and simply store to history. After a certain amount of these joint position discrepancies occur, we add a trajectory waypoint to have the arm move to the correct joint position smoothly.